\section{Appendix: övningar och ordlista}
\label{sec:appendix}

\subsection{Övningar med lösningar}

Övningarna följer kapitlen. De flesta löser du med penna och papper, några genom att köra
motorn. Kommandona förutsätter att du står i repots rot med den virtuella miljön aktiverad.
Alla körs offline och ger därför exakt de tal som står i lösningarna.

\subsubsection*{Redovisning}

\begin{ovning}[penna och papper]
Ett bolag köper en maskin för 100\kr{} kontant den 1 januari och skriver av den linjärt på
fem år. Skattesatsen är 20,6\pct. Hur påverkas de tre rapporterna det första året?
\end{ovning}
\begin{losning}
Avskrivningen är 20. Resultaträkning: resultat före skatt $-20$, skatt $+4,12$, nettoresultat
$-15,88$. Kassaflöde: rörelsen $-15,88 + 20 = +4,12$, investeringar $-100$, kassan $-95,88$.
Balansräkning: anläggningstillgångar $+100 - 20 = +80$, kassa $-95,88$, alltså tillgångar
$-15,88$, och eget kapital $-15,88$. Det går ihop. Lägg märke till att själva köpet inte syns
i resultaträkningen.
\end{losning}

\begin{ovning}[penna och papper]
Ett bolag har EBIT 200, avskrivningar 50, capex 70, en ökning av rörelsekapitalet på 10,
räntekostnader 30 och ökar sina lån med 20 netto. Skattesatsen är 20,6\pct. Beräkna FCFF och
FCFE.
\end{ovning}
\begin{losning}
Enligt~\eqref{eq:fcff}: $\FCFF = 200 \cdot 0,794 + 50 - 70 - 10 = 128,8$.
Enligt~\eqref{eq:fcfe}: $\FCFE = 128,8 - 30 \cdot 0,794 + 20 = 125,0$.
\end{losning}

\subsubsection*{DCF}

\begin{ovning}[penna och papper]
Nästa års kassaflöde är 100, WACC 8\pct{} och evig tillväxt 2\pct. (a) Beräkna slutvärdet. (b)
Hur ändras det om $g$ är 3\pct, respektive om WACC är 9\pct? (c) Hur stor del av värdet kommer
från kassaflöden mer än 20 år bort?
\end{ovning}
\begin{losning}
(a) $100 / 0,06 = 1\,667$. (b) $100 / 0,05 = 2\,000$, alltså $+20$\pct, respektive
$100 / 0,07 = 1\,429$, alltså $-14$\pct. (c) Summan av termerna efter de $n$ första i den
geometriska serien~\eqref{eq:gordon} är andelen $q^n$ av helheten, med
$q = 1,02 / 1,08 = 0,944$. För $n = 20$ blir det $0,944^{20} = 0,32$. En tredjedel av
slutvärdet kommer alltså från år som ligger mer än två decennier efter prognosens slut.
\end{losning}

\begin{ovning}[penna och papper]
Ett bolag ger kassaflödet 100 om året i tre år. Därefter växer kassaflödet med 2\pct{} om året
för alltid, med början på 102 år fyra. WACC är 10\pct. Beräkna rörelsevärdet (a) med
årsslutskonvention och (b) med halvårskonvention. (c) Vad är kvoten mellan svaren, och varför?
\end{ovning}
\begin{losning}
(a) Prognosåren: $100\,(1,1^{-1} + 1,1^{-2} + 1,1^{-3}) = 248,7$. Slutvärde:
$102 / 0,08 = 1\,275$, nuvärde $1\,275 / 1,1^3 = 957,9$. Summa $1\,206,6$. (b) Prognosåren:
$100\,(1,1^{-0,5} + 1,1^{-1,5} + 1,1^{-2,5}) = 260,8$. Slutvärdet får faktorn $1,1^{0,5}$:
$1\,275 \cdot 1,0488 / 1,1^3 = 1\,004,7$. Summa $1\,265,5$. (c) Kvoten är $1,0488 = 1,1^{0,5}$.
Med Gordon-slutvärde flyttar halvårskonventionen \emph{alla} kassaflöden ett halvår närmare,
så hela värdet multipliceras med $(1+r)^{1/2}$. Med exit-multipel gäller det bara prognosåren.
\end{losning}

\begin{ovning}[kör motorn]
Stäng av capex-golvet i slutåret för Atlas Copco. Vad blir Gordon-värdet och den implicita
multipeln? Kontrollera kassaflödet år $N+1$ för hand med~\eqref{eq:terminal-fcff}.
\begin{lstlisting}
from valuation_engine.pipeline import run_valuation
from valuation_engine.models import dcf

d = run_valuation("ATCO-A.ST", offline=True).dcf
out = dcf.dcf_core(**{**d.core_inputs, "terminal_capex_floor": False})
print(out.value_per_share_gordon[0], out.implied_exit_multiple[0], out.terminal_fcff[0] / 1e9)
\end{lstlisting}
\end{ovning}
\begin{losning}
Utskriften är 124,69\kr{} per aktie, multipeln 12,8 och kassaflödet 37,78~mdkr (mot 115,85,
11,7 och 34,38 med golvet). För hand: utan golv är
$u = 0,2135 \cdot 0,794 + 0,0504 - 0,0341 = 0,1859$ och
$\FCFF_{N+1} = 204,7 \cdot (1,02 \cdot 0,1859 - 0,2516 \cdot 0,02) = 37,8$. Golvet kostar
alltså ungefär 9\kr{} per aktie.
\end{losning}

\subsubsection*{WACC och beta}

\begin{ovning}[penna och papper]
Ett bolag har börsvärde 800 och skuld 200. Beta är 1,2, den riskfria räntan 3,075\pct,
riskpremien 4,23\pct, låneräntan 5\pct{} och skattesatsen 20,6\pct. Beräkna WACC.
\end{ovning}
\begin{losning}
$r_e = 3,075 + 1,2 \cdot 4,23 = 8,15$\pct. Låneränta efter skatt: $5 \cdot 0,794 = 3,97$\pct.
$\WACC = 0,8 \cdot 8,15 + 0,2 \cdot 3,97 = 7,31$\pct.
\end{losning}

\begin{ovning}[penna och papper]
Fyra veckors avkastning i procent för marknaden: $2, -1, 3, 0$. För aktien:
$3, -2, 4, -1$. Skatta beta med~\eqref{eq:beta}, Blume-justera, och beräkna $R^2$.
\end{ovning}
\begin{losning}
Båda medelvärdena är 1. Avvikelserna är $(1, -2, 2, -1)$ för marknaden och $(2, -3, 3, -2)$
för aktien. Summan av produkterna är $2 + 6 + 6 + 2 = 16$ och marknadens kvadratsumma
$1 + 4 + 4 + 1 = 10$, så $\hat\beta = 16/10 = 1,6$. (Nämnaren $n - 1$ förkortas bort.)
Blume: $\tfrac23 \cdot 1,6 + \tfrac13 = 1,4$. Aktiens kvadratsumma är 26, så
$R^2 = 16^2 / (10 \cdot 26) = 0,98$. Fyra observationer är naturligtvis alldeles för få. Motorn
kräver 52.
\end{losning}

\begin{ovning}[penna och papper]
Ett jämförelsebolag har beta 1,3 och $D/E = 0,5$. Målbolaget har $D/E = 0,2$. Skattesatsen är
20,6\pct. Beräkna målbolagets beta med~\eqref{eq:hamada} och Blume-justera det, som motorn gör.
\end{ovning}
\begin{losning}
Utan skuld: $\beta_U = 1,3 / (1 + 0,794 \cdot 0,5) = 0,93$. Med målbolagets skuld:
$\beta_L = 0,93 \cdot (1 + 0,794 \cdot 0,2) = 1,08$. Blume:
$\tfrac23 \cdot 1,08 + \tfrac13 = 1,05$.
\end{losning}

\subsubsection*{Multiplar och klustring}

\begin{ovning}[penna och papper]
Fem jämförelsebolag handlas till 8, 9, 10, 12 och 14 gånger EBITDA. Målbolaget har EBITDA 500,
nettoskuld 800 och 100 aktier. (a) Beräkna lågt, bas och högt värde per aktie som motorn gör.
(b) Vilken evig tillväxt motsvarar medianmultipeln enligt~\eqref{eq:multipel-dcf} om
kassakonverteringen är 0,6 och WACC 8\pct?
\end{ovning}
\begin{losning}
(a) 25:e percentilen, medianen och 75:e percentilen är 9, 10 och 12. Rörelsevärde 4\,500,
5\,000 och 6\,000. Minus nettoskuld: 3\,700, 4\,200 och 5\,200. Per aktie: 37, 42 och 52\kr.
(b) $10 = 0,6 / (0,08 - g)$ ger $g = 2$\pct.
\end{losning}

\begin{ovning}[penna och papper]
Fem bolag har z-värdena 1, 2, 3, 8 och 9 i en enda variabel. Kör Lloyds algoritm med $k = 2$
från startpunkterna $\mu_1 = 1$ och $\mu_2 = 4$. Beräkna sedan $J$ och silhuettvärdet för
punkten 3.
\end{ovning}
\begin{losning}
Varv 1: punkterna 1 och 2 går till $\mu_1$, punkterna 3, 8 och 9 till $\mu_2$. Nya
medelpunkter: 1,5 och 6,67. Varv 2: punkten 3 ligger nu närmare 1,5 (avstånd 1,5 mot 3,67) och
byter kluster. Nya medelpunkter: 2 och 8,5. Varv 3: ingen punkt byter, algoritmen är klar med
klustren $\{1, 2, 3\}$ och $\{8, 9\}$. $J = (1 + 0 + 1) + (0,25 + 0,25) = 2,5$. För punkten 3:
$a = (2 + 1)/2 = 1,5$ och $b = (5 + 6)/2 = 5,5$, så $s = (5,5 - 1,5)/5,5 = 0,73$.
Medelvärdet över alla fem punkter är 0,81, en tydlig struktur. Jämför med 0,38 i
figur~\ref{fig:silhuett}.
\end{losning}

\begin{ovning}[kör motorn]
Värdera Sandvik, som också ligger bland fixtures. Vilket $k$ väljs, är Atlas Copco bland
jämförelsebolagen, och varför är exit-multipeln exakt 20,0?
\begin{lstlisting}[style=sh]
python -m valuation_engine SAND.ST --offline --no-charts --no-excel
\end{lstlisting}
\end{ovning}
\begin{losning}
Även här väljs $k = 2$, och Atlas Copco är det närmaste jämförelsebolaget. Comps ger
429,5\kr{} på EV/EBITDA mot kursen 362,10, medan Gordon-DCF:en ger 210,0. Det är samma mönster
som i caset. Jämförelsebolagens median är 20,2 gånger EBITDA, men DCF:ens standardvärde för
exit-multipeln begränsas till intervallet 4 till 20, och utskriften säger det: ”20.2x clipped
to 20.0x”.
\end{losning}

\subsubsection*{LBO och IRR}

\begin{ovning}[penna och papper]
Ett bolag med EBITDA 200 köps för 8 gånger EBITDA med lån på 5 gånger EBITDA (bortse från
avgifter). Efter fem år är EBITDA 260, nettoskulden 600 och bolaget säljs för 9 gånger EBITDA.
Beräkna MOIC och IRR, och dela upp värdeskapandet enligt~\eqref{eq:attribution}.
\end{ovning}
\begin{losning}
Köp: rörelsevärde 1\,600, lån 1\,000, insats 600. Försäljning: $9 \cdot 260 = 2\,340$, minus
skuld 600 ger 1\,740. $\MOIC = 1\,740 / 600 = 2,9$ och $\IRR = 2,9^{1/5} - 1 = 23,7$\pct.
Uppdelning: resultattillväxt $60 \cdot 8 = 480$, multipel $1 \cdot 260 = 260$, amortering
$1\,000 - 600 = 400$. Summa $1\,140 = 1\,740 - 600$, och andelarna är 42, 23 och 35\pct.
\end{losning}

\begin{ovning}[penna och papper]
(a) Beräkna internräntan för kassaflödena $-100, +60, +60$. (b) Visa att flödena
$-100, +230, -132$ har två internräntor. Vilken returnerar motorns \kod{compute_irr}?
\end{ovning}
\begin{losning}
(a) Med $x = 1/(1+r)$: $60x^2 + 60x - 100 = 0$, så
$x = (-60 + \sqrt{3\,600 + 24\,000})/120 = 0,884$ och $r = 13,1$\pct. Den andra roten är
negativ och saknar betydelse. (b) Multiplicera~\eqref{eq:irr} med $(1+r)^2$:
$100(1+r)^2 - 230(1+r) + 132 = 0$, så $1 + r = (230 \pm 10)/200$, alltså $r = 10$\pct{} eller
$r = 20$\pct. Två teckenbyten tillåter två positiva rötter. Motorn returnerar den minsta,
10\pct.
\end{losning}

\begin{ovning}[kör motorn]
Kontrollera motorns betalningsförmåga: kör en LBO på Atlas Copco till det pris som ska ge
20\pct{} och läs av avkastningen. Vilken av de tre källorna dominerar, och vad säger det?
\begin{lstlisting}
from valuation_engine.pipeline import run_valuation
from valuation_engine.models import lbo

r = run_valuation("ATCO-A.ST", offline=True)
price = r.lbo.ability_to_pay[0.20]                       # 117.93
x = lbo.run_lbo(r.company, lbo.with_entry_price(r.lbo.assumptions, price))
print(x.irr, x.moic, x.entry_multiple)
print(x.attribution)
\end{lstlisting}
\end{ovning}
\begin{losning}
IRR är 20,0\pct, MOIC 2,49 och köpmultipeln 13,3. Uppdelningen i mdkr: resultattillväxt 132,9,
multipelförändring 329,1, amortering 98,6 och avgifter $-16,6$, summa 544,0. Multipeln står
för 60\pct{} av värdeskapandet: bolaget köps till 13,3 och säljs till jämförelsebolagens 19,4
gånger EBITDA. LBO-värdet vilar alltså tungt på samma exit-multipel som exit-DCF:en, vilket
är en av casets redovisade svagheter.
\end{losning}

\subsubsection*{Monte Carlo}

\begin{ovning}[penna och papper]
(a) Bestäm Choleskyfaktorn för korrelationsmatrisen med $\rho = 0,3$ i två dimensioner. (b)
Avbilda dragningen $z = (1,0;\; -0,5)$. (c) Visa att $x_2$ har varians 1 och att
$\Cov(x_1, x_2) = \rho$.
\end{ovning}
\begin{losning}
(a) Enligt~\eqref{eq:cholesky2}: $L_{21} = 0,3$ och $L_{22} = \sqrt{1 - 0,09} = 0,954$. (b)
$x_1 = 1,0$ och $x_2 = 0,3 \cdot 1,0 + 0,954 \cdot (-0,5) = -0,177$. (c)
$\Var(x_2) = \rho^2 + (1 - \rho^2) = 1$ och
$\Cov(x_1, x_2) = \Cov(z_1,\, \rho z_1 + \sqrt{1-\rho^2}\, z_2) = \rho$, eftersom $z_1$ och $z_2$
är oberoende.
\end{losning}

\begin{ovning}[penna och papper]
Simuleringen för Atlas Copco har standardavvikelsen 32,7\kr. (a) Hur många dragningar krävs
för ett medelfel på högst 0,10\kr? (b) Uppskatta med Taylorutvecklingen i
avsnitt~\ref{sec:mc} hur mycket medelvärdet lyfts över basfallet av osäkerheten i WACC
ensam. Den trunkerade fördelningens faktiska standardavvikelse är 0,88 procentenheter och
$r - g = 5,65$.
\end{ovning}
\begin{losning}
(a) $n \ge (32,7 / 0,10)^2 \approx 107\,000$. (b) Den relativa höjningen är
$\sigma^2/(r-g)^2 = (0,88/5,65)^2 = 2,4$\pct. Motorn bekräftar: med standardavvikelsen noll
för de tre andra drivarna blir medelvärdet 118,9\kr, vilket är 2,6\pct{} över basfallets
115,8, medan medianen stannar på 116,0.
\end{losning}

\begin{ovning}[kör motorn]
Kör simuleringen med 100\,000 dragningar och slumpfrö 7. Hur mycket flyttar sig percentilerna
jämfört med standardkörningen, och vad är slutsatsen?
\begin{lstlisting}[style=sh]
python -m valuation_engine ATCO-A.ST --offline --no-charts --no-excel --simulations 100000 --seed 7
\end{lstlisting}
\end{ovning}
\begin{losning}
Utskriften ger P10 84,01, P50 116,32 och P90 165,18, mot 84,29, 116,74 och 164,92 med 10\,000
dragningar och slumpfrö 42. Skillnaderna är några tiondels kronor. Simuleringsfelet är alltså
försumbart jämfört med osäkerheten i antagandena, och 10\,000 dragningar räcker.
\end{losning}

\clearpage
\subsection{Ordlista}

Svenska termer i bokstavsordning med de engelska motsvarigheter som används i intervjuer.

\begingroup
\small
\setlength{\columnsep}{1.6em}\setlength{\parskip}{0pt}
\begin{multicols}{2}
\raggedright
\newcommand{\ord}[2]{\noindent\textbf{#1}\quad\textenglish{\emph{#2}}\par\vspace{0.22em}}
\ord{aktievärde, börsvärde}{equity value, market capitalization}
\ord{amortering}{debt repayment, amortization}
\ord{anläggningstillgångar}{property, plant and equipment (PP\&E), fixed assets}
\ord{avkastningskrav på eget kapital}{cost of equity}
\ord{avskrivningar}{depreciation and amortization (D\&A)}
\ord{balansräkning}{balance sheet}
\ord{betalningsförmåga}{ability to pay}
\ord{brygga}{equity bridge, EV to equity bridge}
\ord{bruttoresultat}{gross profit}
\ord{budpremie}{takeover premium}
\ord{börsintroduktion}{initial public offering (IPO)}
\ord{checkkredit}{revolving credit facility (revolver)}
\ord{diskonteringsränta}{discount rate}
\ord{efterställt lån}{subordinated debt, mezzanine}
\ord{eget kapital}{shareholders' equity}
\ord{evig tillväxt}{terminal growth, perpetual growth rate}
\ord{fritt kassaflöde}{free cash flow (FCF)}
\ord{förvärv}{acquisition}
\ord{halvårskonvention}{mid-year convention}
\ord{hävstång}{leverage}
\ord{internränta}{internal rate of return (IRR)}
\ord{investeringar}{capital expenditures (capex)}
\ord{jämförelsebolag}{comparable companies, peers}
\ord{jämförbara transaktioner}{precedent transactions}
\ord{kassaflödesanalys}{cash flow statement}
\ord{kassakonvertering}{cash conversion}
\ord{kundfordringar}{accounts receivable}
\ord{källor och användning}{sources and uses}
\ord{känslighetsanalys}{sensitivity analysis}
\ord{lager}{inventory}
\ord{leverantörsskulder}{accounts payable}
\ord{låneränta}{cost of debt}
\ord{medelfel}{standard error}
\ord{minoritetsintressen}{minority interest, non-controlling interests}
\ord{multipel}{multiple}
\ord{nettoresultat}{net income}
\ord{nettoskuld}{net debt}
\ord{nuvärde}{present value (PV)}
\ord{nyemission}{rights issue, share issue}
\ord{omsättning}{revenue, sales}
\ord{resultaträkning}{income statement}
\ord{riskfri ränta}{risk-free rate}
\ord{riskkapitalbolag}{private equity firm, financial sponsor}
\ord{riskpremie}{equity risk premium (ERP)}
\ord{rullande tolv månader}{last twelve months (LTM)}
\ord{rörelsekapital}{net working capital (NWC)}
\ord{rörelseresultat}{operating income, EBIT}
\ord{rörelsevärde}{enterprise value (EV)}
\ord{seniort lån}{senior debt, term loan}
\ord{skattesköld}{tax shield}
\ord{skuldsättning}{leverage, gearing}
\ord{slutvärde}{terminal value (TV)}
\ord{storlekspremie}{size premium}
\ord{stubbperiod}{stub period}
\ord{uppköp, lånefinansierat}{leveraged buyout (LBO)}
\ord{utdelning}{dividend}
\ord{vägd kapitalkostnad}{weighted average cost of capital (WACC)}
\end{multicols}
\endgroup
